\begin{pastexam}[label={px:gw3-101-ds-03}]{關務三等資訊處理 101 年 第三題}
（申論）

當飛機航班客滿或超額預約（overbooked）時，在機場航空公司櫃台常有機位候補（standby）的長龍。候補優先次序是由
乘客的購票價格、累積里程數及要求候補的時間先後等因素決定優先權（priority）。請設計一個有效率的資料結構及其運算
（operations）來處理候補，並說明其時間複雜度。（15 分）
\end{pastexam}

\begin{pastsoln}
答案：用以「優先權」為鍵的\textbf{最大堆積}實作優先佇列（資結書第~8.1、8.3~節）。 加入候補、取出下一位都是
$O(\log n)$，查看下一位是 $O(1)$（$n$ 是候補人數）。

\textbf{優先權怎麼比。} 題目沒有規定三個因素的權重，這裡採字典順序：先比購票價格（高者優先），相同再比累積里程
（多者優先），再相同就比要求候補的時間（早者優先）。 每位乘客登記時給一個遞增的序號 $\mathit{seq}$ 代表時間先後，
所以任兩位乘客的優先權一定分得出高低，同條件的人也照先來後到（公平，相當於穩定）。 定義
\[
a \succ b \iff (a.\mathit{price}, a.\mathit{miles}, -a.\mathit{seq}) > (b.\mathit{price}, b.\mathit{miles}, -b.\mathit{seq})
\quad\text{（字典順序）}.
\]
若航空公司改用加權分數（例如 $w_1 \cdot$ 價格 $+ w_2 \cdot$ 里程），只要換掉這個比較函式，結構不變。

\textbf{資料結構。}
\begin{itemize}
\item 乘客紀錄：姓名或編號 $\mathit{id}$、$\mathit{price}$、$\mathit{miles}$、$\mathit{seq}$。
\item 最大堆積 $H[1..n]$（陣列存的完整二元樹），以 $\succ$ 比較：每個節點的優先權不低於子節點，根就是下一位該
補上的乘客。
\item 索引表 $\mathit{pos}$：由 $\mathit{id}$ 查到乘客目前在 $H$ 的哪個位置（雜湊表），供取消候補用（帶索引的堆積，
資結書第~8.3~節）。
\end{itemize}

\textbf{運算與時間複雜度。}
\begin{center}
\begin{tabular}{p{0.2\linewidth}p{0.5\linewidth}l}
\toprule
運算 & 做法 & 時間 \\
\midrule
登記候補 & 放在 $H$ 最後，往上和父節點比、優先權較高就交換（上浮），同時更新 $\mathit{pos}$ & $O(\log n)$ \\
查看下一位 & 讀 $H[1]$ & $O(1)$ \\
補上一位（有空位時） & 取出 $H[1]$，把最後一個搬到根，往下和優先權較高的子節點交換（下沉） & $O(\log n)$ \\
取消候補 & 由 $\mathit{pos}$ 找到位置，和最後一個交換後刪掉，再視情況上浮或下沉 & $O(\log n)$ \\
資格變更（例如升等） & 改鍵值後上浮或下沉 & $O(\log n)$ \\
\bottomrule
\end{tabular}
\end{center}
一開始已有 $n$ 位乘客時，可以用由下而上建堆積，$O(n)$（資結書第~8.2~節）。 開放 $k$ 個空位時，取出 $k$ 次，
$O(k \log n)$。

\textbf{為什麼不用別的。} 依序排好的陣列或串列，查看、取出是 $O(1)$，但每有一人登記就要插入到正確位置，$O(n)$；
不排序的串列，登記 $O(1)$ 但每次找最高優先權要 $O(n)$。 候補時登記和補位都很頻繁，堆積兩種操作都是
$O(\log n)$，最平衡。 （平衡搜尋樹也可以全部做到 $O(\log n)$，而且能依序列出整份候補名單，但實作比較複雜。）
\end{pastsoln}
